\documentclass[10pt,a4paper]{article}

\usepackage[a4paper,top=13mm,bottom=15mm,left=15mm,right=15mm]{geometry}
\usepackage{fontspec}
\usepackage[table]{xcolor}
\usepackage{tabularx}
\usepackage{array}
\usepackage[most]{tcolorbox}
\usepackage{fancyhdr}
\usepackage{lastpage}
\usepackage{microtype}
\usepackage{ragged2e}

\IfFontExistsTF{Arial}{
  \setmainfont{Arial}
  \setsansfont{Arial}
}{
  \setmainfont{TeX Gyre Heros}
  \setsansfont{TeX Gyre Heros}
}

\definecolor{AIEPRed}{HTML}{D62027}
\definecolor{AIEPRedSoft}{HTML}{FDEBEC}
\definecolor{AIEPNavy}{HTML}{082B4C}
\definecolor{AIEPNavyMid}{HTML}{17476F}
\definecolor{AIEPBlueSoft}{HTML}{EDF4FA}
\definecolor{AIEPInk}{HTML}{1F2933}
\definecolor{AIEPMuted}{HTML}{536171}
\definecolor{AIEPLine}{HTML}{D8E1EA}
\definecolor{AIEPSoft}{HTML}{F6F8FA}
\definecolor{AIEPGold}{HTML}{C59A32}
\definecolor{AIEPGoldSoft}{HTML}{FBF5E6}
\definecolor{AIEPGreen}{HTML}{26734D}
\definecolor{AIEPGreenSoft}{HTML}{EAF5EF}

\pagestyle{fancy}
\fancyhf{}
\renewcommand{\headrulewidth}{0pt}
\fancyfoot[L]{\footnotesize\textcolor{AIEPMuted}{GeoGreen Escolar · AIEP Osorno}}
\fancyfoot[C]{\footnotesize\textcolor{AIEPMuted}{Talleres y mentorías · 2026}}
\fancyfoot[R]{\footnotesize\textcolor{AIEPMuted}{\thepage\ de \pageref{LastPage}}}

\setlength{\parindent}{0pt}
\setlength{\parskip}{4pt}
\setlength{\tabcolsep}{5pt}
\RaggedRight
\hyphenpenalty=8000
\exhyphenpenalty=8000
\sloppy

\newcolumntype{Y}{>{\RaggedRight\arraybackslash}X}
\newcolumntype{L}[1]{>{\RaggedRight\arraybackslash}p{#1}}

\newtcolorbox{hero}{
  enhanced,
  colback=AIEPNavy,
  colframe=AIEPNavy,
  boxrule=0pt,
  arc=2mm,
  left=8mm,right=8mm,top=6mm,bottom=5.5mm,
  before skip=0pt,after skip=5mm,
  borderline south={1.2mm}{0pt}{AIEPRed}
}

\newtcolorbox{modebox}[3][]{
  enhanced,
  colback=#2,
  colframe=#3,
  boxrule=0.7pt,
  arc=1.8mm,
  left=4.5mm,right=4.5mm,top=4mm,bottom=4mm,
  before skip=0pt,after skip=0pt,
  #1
}

\newtcolorbox{callout}[2][]{
  enhanced,
  colback=AIEPSoft,
  colframe=#2,
  boxrule=0.7pt,
  arc=1.6mm,
  left=4.5mm,right=4.5mm,top=3mm,bottom=3mm,
  before skip=4mm,after skip=0pt,
  #1
}

\newtcolorbox{sectionbar}[2][]{
  enhanced,
  colback=#2,
  colframe=#2,
  boxrule=0pt,
  arc=1.2mm,
  left=4mm,right=4mm,top=2.4mm,bottom=2.4mm,
  before skip=4.5mm,after skip=2mm,
  #1
}

\newcommand{\metric}[2]{%
  {\fontsize{23}{25}\selectfont\bfseries\color{#1}#2\par}%
}
\newcommand{\ownerred}[1]{{\scriptsize\bfseries\color{AIEPRed}\MakeUppercase{#1}}}
\newcommand{\ownernavy}[1]{{\scriptsize\bfseries\color{AIEPNavyMid}\MakeUppercase{#1}}}
\newcommand{\deliverable}[1]{\par\vspace{1mm}{\small\textcolor{AIEPMuted}{\textbf{Producto:} #1}}}

\begin{document}
\fontsize{10.4}{13.2}\selectfont

\begin{hero}
{\small\bfseries\color{white!75}\MakeUppercase{Propuesta de ejecución · 2026}\par}
\vspace{3mm}
{\fontsize{27}{31}\selectfont\bfseries\color{white}GeoGreen Escolar\par}
\vspace{1mm}
{\Large\color{white!84}Organización, propósitos y productos\par}
\vspace{4mm}
{\small\color{white!70}Instituto Comercial Liceo Bicentenario · Osorno\par}
\end{hero}

\begin{callout}{AIEPNavyMid}
Cuatro talleres construyen una base común. Cuatro mentorías acompañan a los
equipos que quieran orientación para desarrollar y comunicar sus propuestas.
El recorrido culmina en un concurso final con jurado y premiación.
\end{callout}

\vspace{5mm}
{\Large\bfseries\color{AIEPNavy}La jornada, de un vistazo\par}
\vspace{2.5mm}
\begin{tabularx}{\textwidth}{@{}Y@{\hspace{3.5mm}}Y@{}}
\begin{modebox}[height=48mm]{AIEPRedSoft}{AIEPRed}
{\scriptsize\bfseries\color{AIEPMuted}TALLERES · COHORTE COMPLETA\par}
\vspace{2mm}
\metric{AIEPRed}{2 × 90 min}
\textbf{30 estudiantes y 5 equipos por bloque}\par
\vspace{2mm}
El mismo taller se realiza con ambos bloques.\par
\textbf{Cobertura:} aproximadamente 60 estudiantes.\par
\textbf{Duración total diaria:} 3 horas.
\end{modebox}
&
\begin{modebox}[height=48mm]{AIEPBlueSoft}{AIEPNavyMid}
{\scriptsize\bfseries\color{AIEPMuted}MENTORÍAS · APOYO AL DESARROLLO\par}
\vspace{2mm}
\metric{AIEPNavy}{2 × 90 min}
\textbf{Para los equipos que quieran apoyo}\par
\vspace{2mm}
Permiten revisar decisiones, contrastar evidencia y destrabar un próximo paso.\par
\textbf{La propuesta sigue siendo desarrollada por el equipo.}
\end{modebox}
\end{tabularx}

\begin{sectionbar}{AIEPRed}
{\bfseries\color{white}CUATRO TALLERES · UNA BASE COMÚN · 90 MINUTOS POR BLOQUE}
\end{sectionbar}

\renewcommand{\arraystretch}{1.28}
\begin{tabularx}{\textwidth}{@{}L{25mm}L{45mm}Y@{}}
\rowcolor{AIEPRedSoft}
\textbf{17 AGO} &
\textbf{Taller 1}\par\ownerred{Desarrollo Social}\par Conciencia ambiental &
Observar residuos, hábitos y lugares para formular una situación ambiental concreta.
\deliverable{Equipos organizados y problema definido.} \\[2.2mm]

\rowcolor{white}
\textbf{18 AGO} &
\textbf{Taller 2}\par\ownerred{Desarrollo Social}\par Ciencia del reciclaje &
Comprender el residuo o material, su reciclabilidad y las condiciones para separarlo correctamente.
\deliverable{Ficha de análisis del residuo o material.} \\[2.2mm]

\rowcolor{AIEPBlueSoft}
\textbf{24 AGO} &
\textbf{Taller 3}\par\ownernavy{I/E/T}\par Sensores y GeoGreen &
Relacionar el problema con una variable, un sensor y una primera prueba segura.
\deliverable{Idea tecnológica inicial y plan de prueba.} \\[2.2mm]

\rowcolor{white}
\textbf{25 AGO} &
\textbf{Taller 4}\par\ownernavy{I/E/T}\par Del dato a la acción &
Usar datos simulados para definir estados, una interfaz, una alerta y una decisión.
\deliverable{Prototipo digital provisional y evidencia de prueba.} \\
\end{tabularx}
\renewcommand{\arraystretch}{1}

\begin{callout}{AIEPRed}
\textbf{RESULTADO DE LOS TALLERES}\quad
Problema concreto $\rightarrow$ residuo o material $\rightarrow$ idea tecnológica
$\rightarrow$ primera representación de software útil.
\end{callout}

\newpage

{\small\bfseries\color{AIEPGold}\MakeUppercase{GeoGreen Escolar · Acompañamiento a los equipos}\par}
\vspace{1.5mm}
{\fontsize{23}{27}\selectfont\bfseries\color{AIEPNavy}Cuatro mentorías con un propósito preciso\par}
\vspace{1mm}
\textcolor{AIEPMuted}{Dos bloques de 90 minutos por fecha, disponibles para los equipos que quieran orientación.}

\begin{sectionbar}{AIEPGold}
{\bfseries\color{white}REVISAR · ORIENTAR · DESTRABAR · DEFINIR EL PRÓXIMO PASO}
\end{sectionbar}

\renewcommand{\arraystretch}{1.3}
\begin{tabularx}{\textwidth}{@{}L{25mm}L{45mm}Y@{}}
\rowcolor{AIEPRedSoft}
\textbf{31 AGO} &
\textbf{Mentoría 1}\par\ownerred{Desarrollo Social}\par Problema y contexto &
Contrastar el problema con el contexto y las personas relacionadas; distinguir hechos de suposiciones.
\deliverable{Problema validado, contexto afectado y próximo paso.} \\[2.5mm]

\rowcolor{white}
\textbf{7 SEP} &
\textbf{Mentoría 2}\par\ownernavy{I/E/T}\par Solución y recursos &
Ajustar la solución mínima, los componentes, el software, las tareas y los recursos necesarios.
\deliverable{Solución mínima y recursos definidos.} \\[2.5mm]

\rowcolor{AIEPBlueSoft}
\textbf{21 SEP} &
\textbf{Mentoría 3}\par\ownernavy{I/E/T}\par Avance y evidencia &
Revisar una maqueta, simulación, diagrama, prototipo o prueba y priorizar la siguiente mejora.
\deliverable{Evidencia revisada y ajuste priorizado.} \\[2.5mm]

\rowcolor{white}
\textbf{28 SEP} &
\textbf{Mentoría 4}\par\ownerred{Desarrollo Social}\par Comunicación y ensayo &
Ordenar el relato, preparar el apoyo visual, distribuir las intervenciones y ensayar la presentación.
\deliverable{Guion, soporte visual, roles, tiempos y presentación corregida.} \\
\end{tabularx}
\renewcommand{\arraystretch}{1}

\begin{sectionbar}{AIEPRed}
{\bfseries\color{white}SEMANA DEL 14 DE SEPTIEMBRE · SIN ACTIVIDADES · FIESTAS PATRIAS}
\end{sectionbar}

\vspace{4mm}
{\Large\bfseries\color{AIEPNavy}Responsables\par}
\vspace{2mm}
\renewcommand{\arraystretch}{1.32}
\begin{tabularx}{\textwidth}{@{}L{58mm}Y@{}}
\rowcolor{AIEPRedSoft}
{\bfseries\color{AIEPRed}DESARROLLO SOCIAL} & Talleres 1 y 2 · Mentorías 1 y 4 \\[1.3mm]
\rowcolor{AIEPBlueSoft}
{\bfseries\color{AIEPNavyMid}INGENIERÍA, ENERGÍA Y TECNOLOGÍA (I/E/T)} & Talleres 3 y 4 · Mentorías 2 y 3 \\[1.3mm]
\rowcolor{AIEPGreenSoft}
{\bfseries\color{AIEPGreen}EQUIPO CONJUNTO} & Hito comunicacional · Evento final · Cierre interno \\
\end{tabularx}
\renewcommand{\arraystretch}{1}

\vspace{5mm}
{\Large\bfseries\color{AIEPNavy}Aspectos de coordinación\par}
\vspace{2mm}
\begin{tabularx}{\textwidth}{@{}Y@{\hspace{4mm}}Y@{}}
\begin{modebox}[height=31mm]{AIEPSoft}{AIEPLine}
\textbf{Para confirmar con el establecimiento}\par
\vspace{1.5mm}
Fechas definitivas · cantidad específica de estudiantes · espacios para ambos bloques.
\end{modebox}
&
\begin{modebox}[height=31mm]{AIEPSoft}{AIEPLine}
\textbf{Para organizar las jornadas}\par
\vspace{1.5mm}
Horarios de inicio y término · posibilidad de realizar mentorías en el liceo y en AIEP.
\end{modebox}
\end{tabularx}

\begin{sectionbar}{AIEPGreen}
{\bfseries\color{white}CONTINUIDAD · 2 OCT HITO COMUNICACIONAL · 5 OCT CONCURSO FINAL, JURADO Y PREMIACIÓN · SEMANA 12 OCT CIERRE INTERNO}
\end{sectionbar}

\end{document}
